\chap{结论与展望}

\sect{主要结论}

本文研究了金属热处理温度场的受控递归建模问题。模型以完整离散温度场为状态，以炉温为动作，并显式接收边界与材料参数。C45/AISI 1045 类中碳钢数值环境包含温变导热系数、温变比热和对流辐射边界。主要结论如下。

\begin{enumerate}
  \item 完整轨迹评价揭示了一步误差无法表达的长期行为。在相同数据、训练轮数和随机种子下，五步展开相对一步训练将 ID 轨迹 RMSE 降低 $4.2\%$，但控制 OOD 轨迹 RMSE 增加 $10.8\%$。因此，本文依据验证集轨迹误差选择五步设置，并把完整闭环轨迹作为模型选择与报告单位。
  \item 物理约束的主要价值体现在能量一致性和参数组合外推。验证选择的 $10^{-3}$ 权重将控制 OOD 物理残差降低 $44.8\%$。控制 OOD 的总体 RMSE 均值降低 $6.8\%$，但逐轨迹配对区间跨越零，所以该数值只表示当前样本上的趋势。四参数组合 OOD 的滚动 RMSE 降低 $29.5\%$，最大绝对误差降低约 $49.3\%$。
  \item 在本文逐时自由变化的瞬时观测模型下，动态对流辐射边界只能识别 $\Heff=h+\varepsilon h_{\mathrm{rad}}$。该参数化将动态边界滚动 RMSE 从 $2.285\pm0.108$~\si{\degreeCelsius} 降至 $1.409\pm0.068$~\si{\degreeCelsius}。交叉求解器数值验证保持相同排序。
  \item 无中心测点的近表面四点方案取得 \SI{2.128}{\degreeCelsius} 的未观测节点 RMSE 和 \SI{2.854}{\degreeCelsius} 的中心 RMSE。包含中心热电偶的理想五点方案将未观测节点 RMSE 降至 \SI{1.073}{\degreeCelsius}。ID 条件下后验均值控制 30/30 个回合成功。严格参数 OOD 条件下，风险感知控制的综合目标平均配对差为 $-1.737$，95\% bootstrap 区间为 $[-3.291,-0.226]$。
  \item $\Heff$ 世界模型在 BDF 交叉求解器对照中获得 55.4 倍完整轨迹推演加速，在常值尾部候选动作规划中获得 63.4 倍加速。两项结果分别对应状态轨迹生成和滚动决策，说明该模型可以为高保真复核前的批量工艺筛选提供计算余量。
  \item 固定三维试块实验建立了 $9\times7\times5$ 基础对照，并将模型网格进一步加密至 $17\times13\times9$。加密后共有 1989 个温度状态，数据模型和物理模型的 ID 滚动 RMSE 分别为 $3.983$ 和 \SI{4.082}{\degreeCelsius}，相对基础网格降低 $10.8\%$ 和 $11.9\%$。代表时刻的八分体剖视图显示，模型能够恢复六面受热条件下的表面至芯部温度梯度。
\end{enumerate}

上述结果支持一个明确判断。热处理世界模型的关键价值在于形成可递归、可观测、可用于动作比较的受控动力学表示。物理损失增强 OOD 稳定性，可辨识参数化减少动态边界歧义，状态估计器负责连接稀疏测量。三者共同决定模型能否进入闭环。

\sect{工程应用路径}

本文方法适合作为现有热处理控制系统的上层预测与决策模块。该分层结构与工业连续加热炉的模型化轨迹规划和预测控制系统一致\cite{steinboeck2011trajectory,steinboeck2013nmpc,yi2017tracking}。现场测量层提供炉温、表面或近表面热电偶数据。集合滤波器据此更新工件温度场和边界参数后验，并将后验成员转换为候选炉温下的有效换热输入。世界模型在每个决策时刻推演多条候选轨迹，计算心部到温误差、截面温差、表面超温和动作变化。选定炉温作为设定值交给底层 PLC 或 PID 回路执行。设备原有的温度上限、功率限制、急停和联锁逻辑保持独立。

\begin{table}[H]
  \centering
  \caption{世界模型面向热处理生产的功能对应关系}
  \label{tab:industrial_application}
  \small
  \begin{tabularx}{\textwidth}{>{\raggedright\arraybackslash}p{2.1cm}>{\raggedright\arraybackslash}p{3.1cm}>{\raggedright\arraybackslash}X>{\raggedright\arraybackslash}p{3.0cm}}
    \toprule
    生产环节 & 在线输入 & 模型输出 & 可支持的决策 \\
    \midrule
    工件状态监测 & 炉温与近表面测温 & 完整温度场、心部温度及其区间 & 虚拟测温与异常温差识别 \\
    到温与保温判定 & 当前状态、目标温度和允许温差 & 未来温度轨迹与预计到温时刻 & 减少过早出炉和无效保温 \\
    边界变化补偿 & 历史温度响应与炉温动作 & $\Heff$ 或 $h,\varepsilon$ 后验分布 & 适应装炉、气流和表面状态变化 \\
    炉温设定优化 & 候选炉温、温度约束和能耗项 & 各候选动作的轨迹代价与风险分位 & 向底层控制器下发炉温设定值 \\
    工艺离线筛选 & 初始状态、材料参数和工艺曲线库 & 批量推演结果及高风险候选 & 缩小高保真仿真和试制范围 \\
    \bottomrule
  \end{tabularx}
\end{table}

上述功能由本文实验中的不同证据支撑。无中心测点方案取得 \SI{2.128}{\degreeCelsius} 的未观测节点 RMSE，可用于表面测温条件下的内部状态估计。完整轨迹推演和候选动作规划分别获得 55.4 倍和 63.4 倍加速，为决策周期内比较多个炉温设定提供了计算余量。严格参数 OOD 实验中，风险分位控制降低了场景配对综合目标，说明边界后验可以直接进入约束决策。$17\times13\times9$ 三维实验进一步表明，同一状态转移形式能够处理包含 1989 个温度状态的固定几何试块。

工程部署可按离线标定、影子运行和受监督闭环三个阶段推进。离线阶段利用仪器化试块和历史炉次校正物性、边界范围及传感器响应。影子运行阶段只记录模型建议，并与热电偶、硬度检验和高保真计算结果对照。模型通过预定验收指标后，再由上层优化器周期性给出炉温设定值。该路径符合数字孪生由模型、在线数据和服务共同更新的工程结构\cite{tao2018digitaltwin,stark2019digitaltwin}，同时保留现有设备控制与安全体系。

\sect{研究范围}

本文证据主线来自一维对称平板数值实验，并包含固定长方体几何的三维扩展。研究对象是金属热处理中的受控热动力学子问题，不覆盖完整的组织性能控制。BDF 参考环境改变了离散网格和时间积分方法，但连续传热方程与材料物性函数保持一致。当前状态不含相变组织、热应力和变形。\SI{350}{\degreeCelsius} 目标跟踪用于检验控制链路，不构成完整的 C45 工业热处理制度。三维扩展采用单随机种子；网格加密实验重新生成了同一组随机工况，但没有训练跨网格共享模型。闭环实验尚未接入真实炉温、热电偶和执行器数据。因此，文中精度用于说明方法对当前热动力学环境的逼近和控制能力。

\sect{后续工作}

后续研究可沿三个方向推进。第一，引入公开热电偶曲线或自建加热实验，对材料物性、表面发射率和传感器响应进行联合标定。第二，将状态模型扩展到二维截面和复杂几何，并比较局部卷积结构、物理信息算子和降阶模型\cite{wang2021pideeponet,derouiche2021induction,gong2026heatrom}。第三，建立材料 OOD 检测和混合规划器。世界模型负责快速筛选，高保真有限元模型复核少量高风险候选动作。该结构能够保留在线速度，也能提高复杂工况中的控制可信度。
